\documentclass[11pt]{article}
\usepackage{amsmath, amssymb, geometry, mathtools, booktabs, longtable}
\usepackage{hyperref}
\geometry{margin=1in}

\title{RG Checkerboard Conditional Flow for 2D $\phi^4$: Code-Level Description and CLI Reference}
\author{Kostas Orginos}
\date{\today}

\begin{document}
\maketitle

\section{Purpose}

This note documents the new checkerboard RG flow implemented in
\texttt{jaxqft/models/phi4\_rg\_checkerboard\_flow.py}, together with its
trainer
\texttt{scripts/phi4/train\_rg\_checkerboard\_flow.py}
and checker
\texttt{scripts/phi4/check\_rg\_checkerboard\_flow.py}.

This path is a new evolution branch. It leaves both
\texttt{jaxqft/models/phi4\_mg.py} and
\texttt{jaxqft/models/phi4\_rg\_cond\_flow.py} unchanged.

\section{High-Level Design}

The model is built from three ideas:
\begin{itemize}
\item recursive $2\times 2$ RG splitting into one coarse mode and three local
fluctuation modes,
\item red/black conditional affine updates acting only on the fluctuation
variables $\eta$,
\item a shifted second blocking obtained by rolling the fine lattice by
$(1,1)$ before repeating the red/black update.
\end{itemize}

At each non-terminal RG level, the model applies a full-field invertible map on
the current lattice. Only after this level map is applied is the field split
into the next coarse field and the residual fine modes that are stored for the
recursive reconstruction.

\section{Lattice, Prior, and RG Split}

Let the fine lattice be $\Lambda_0$ with shape
\[
L\times L,\qquad L=2^K.
\]
The model acts on batched scalar fields
\[
\phi \in \mathbb{R}^{B\times L\times L}.
\]

The latent prior is sitewise standard normal:
\begin{equation}
p_Z(z)=\prod_{x\in\Lambda_0}\frac{1}{\sqrt{2\pi}}
\exp\!\left(-\frac{z_x^2}{2}\right).
\end{equation}

At level $k$ the lattice size is
\[
L_k = L/2^k.
\]
For each $2\times 2$ block we pack the site values into
\[
v=(v_{00},v_{01},v_{10},v_{11})\in\mathbb{R}^4.
\]

\subsection{Average Blocking}

For \texttt{rg\_type="average"}, the code uses the orthonormal basis
\begin{equation}
H=
\begin{pmatrix}
\tfrac12 & \tfrac12 & \tfrac12 & \tfrac12 \\
\tfrac12 & -\tfrac12 & \tfrac12 & -\tfrac12 \\
\tfrac12 & \tfrac12 & -\tfrac12 & -\tfrac12 \\
\tfrac12 & -\tfrac12 & -\tfrac12 & \tfrac12
\end{pmatrix}
\end{equation}
and writes
\begin{equation}
(c,\eta_1,\eta_2,\eta_3)=H\,v.
\end{equation}
Here $c$ is the normalized coarse mode and
\[
\eta=(\eta_1,\eta_2,\eta_3)\in\mathbb{R}^3
\]
spans the fluctuation subspace.

The arithmetic block average is
\[
\bar v = \frac14 \sum_{i\in \text{block}} v_i = \frac12\,c.
\]
This is the coarse quantity used in the conditioner.

\subsection{Select Blocking}

For \texttt{rg\_type="select"}, the code uses
\begin{equation}
c=v_{00},\qquad \eta=(v_{01},v_{10},v_{11}).
\end{equation}

\subsection{Inverse Reconstruction}

At each block the inverse reconstruction is
\[
v = R(c,\eta).
\]
For average blocking this is $v=H^\top(c,\eta)$.

\section{Recursive Structure}

The recursion stops at the last $2\times 2$ lattice. The terminal block is
mapped by an unconditional $4$-component RealNVP.

For each non-terminal level:
\begin{enumerate}
\item apply a level bijection $T_k$ on the full current lattice,
\item split the result into a next coarse field and level-$k$ fine modes,
\item recurse on the coarse field,
\item reconstruct the fine lattice using the stored fine modes.
\end{enumerate}

Thus the model differs from the previous conditional-flow branch: the
checkerboard updates happen on the \emph{full field at a given scale}, not only
on already-extracted fine variables.

\section{Checkerboard Level Map}

Fix one level and suppress the level index. Let
\[
S(x) = (c,\eta)
\]
be the block split of the current lattice field, where $c$ lives on the coarse
block lattice and $\eta\in\mathbb{R}^3$ is attached to each block.

Let the block lattice be decomposed into red and black checkerboards:
\[
\Lambda_{\mathrm{blk}} = R \cup B,
\qquad
R\cap B = \varnothing.
\]
The code uses the standard parity
\[
\chi(y,x) = (y+x)\bmod 2.
\]

\subsection{One Red Pass}

Define masks $M_R$ and $M_B$ that act on block fields by keeping only red or
black sites, respectively. The red update keeps the black variables fixed and
updates only the red variables:
\begin{align}
u_R &= \bigl[M_B\eta,\ \tilde c\bigr], \\
\eta_R' &= \eta_R \odot \exp\bigl(s_R(u_R)\bigr) + t_R(u_R), \\
\eta_B' &= \eta_B.
\end{align}
Here $\tilde c$ is the conditioning coarse field:
\[
\tilde c =
\begin{cases}
\tfrac12 c,& \text{average blocking,} \\
c,& \text{select blocking.}
\end{cases}
\]

\subsection{One Black Pass}

The black update uses the already-updated red variables as conditioning data:
\begin{align}
u_B &= \bigl[M_R\eta',\ \tilde c\bigr], \\
\eta_B'' &= \eta_B' \odot \exp\bigl(s_B(u_B)\bigr) + t_B(u_B), \\
\eta_R'' &= \eta_R'.
\end{align}

\subsection{Shifted Blocking}

After the red/black pass, the field is merged back to the fine lattice,
\[
x' = R(c,\eta'').
\]
Then the fine lattice is rolled by one in both directions:
\[
\rho(x')_{n_1,n_2} = x'_{n_1+1,n_2+1},
\]
with periodic wrapping. The same red/black update is then applied on the
\emph{shifted blocking}, and the fine lattice is rolled back.

This is the mechanism that couples the original blocking to the complementary
$2\times 2$ blocking.

\subsection{One Full Cycle}

One cycle of the level map is
\begin{equation}
T = \rho^{-1}\circ P_{\mathrm{shift}}\circ \rho \circ P_{\mathrm{orig}},
\end{equation}
where $P_{\mathrm{orig}}$ is the red-then-black pass on the original blocking
and $P_{\mathrm{shift}}$ is the analogous pass on the shifted blocking.

The code repeats this cycle \texttt{n\_cycles} times per level.

\section{Why the Map Remains Invertible}

The critical condition is:
\begin{quote}
when updating a set of variables, the conditioner must depend only on variables
that are held fixed during that substep.
\end{quote}

This condition is enforced in the implementation by masking the \emph{target
color everywhere} before the transformer sees the input.

For example, in the red update the transformer input is
\[
u_R = [M_B\eta,\tilde c].
\]
So even if the transformer uses spatial attention over a neighborhood, it can
see only:
\begin{itemize}
\item black $\eta$ variables,
\item coarse variables $c$.
\end{itemize}
It cannot see any red $\eta$ variables, because they are zeroed everywhere in
the conditioning field. Therefore the red update is triangular with respect to
the red variables.

The same statement holds for the black update, with the roles of red and black
reversed.

\subsection{Does Larger Attention Radius Break This?}

No, \emph{provided the masking is applied before attention}, as in the current
code.

Increasing \texttt{attn\_radius} lets the conditioner use a wider neighborhood
of \emph{frozen} variables. It does \emph{not} introduce dependence on the
variables being updated, because those variables are masked out globally in the
conditioner input.

So:
\begin{itemize}
\item larger \texttt{attn\_radius} increases cost,
\item larger \texttt{attn\_radius} increases the frozen conditioning context,
\item larger \texttt{attn\_radius} does \emph{not} break invertibility in the
current implementation.
\end{itemize}

If one were to feed the unmasked $\eta$ field into the transformer, then yes,
larger attention radius would generally destroy the triangular structure and the
map would no longer be an exact RealNVP-style coupling.

\section{Local Conditioner}

The checkerboard branch now supports
\[
\texttt{conditioner}\in\{\texttt{transformer},\texttt{mlp}\}.
\]

\subsection{MLP Conditioner}

For \texttt{conditioner="mlp"}, each red or black subupdate uses a local
pointwise MLP
\[
(s,t)=\mathrm{MLP}(u_b)
\]
applied independently at each block site $b$. There is no spatial attention in
this case. The conditioning context is still the same masked field
$u_b=[M_{\mathrm{other}}\eta,\tilde c]$, so invertibility is preserved for the
same reason as in the transformer case.

\subsection{Transformer Conditioner}

For \texttt{conditioner="transformer"}, each conditioner is a local transformer
on the block lattice.

\subsection{Token Features}

At each block location $b$ the token is
\[
u_b \in \mathbb{R}^4,
\]
consisting of:
\begin{itemize}
\item the masked conditioning fluctuation vector in $\mathbb{R}^3$,
\item the local coarse scalar.
\end{itemize}

\subsection{Attention Radius}

For \texttt{attn\_radius = r}, the transformer at block $b$ attends to the
window
\[
\mathcal N_r(b)=\{b+\delta:\delta_y,\delta_x\in\{-r,\dots,r\}\}.
\]
Hence:
\begin{itemize}
\item \texttt{attn\_radius = 0} means purely on-site conditioning,
\item \texttt{attn\_radius = 1} means a $3\times 3$ neighborhood on the block
lattice,
\item in general the receptive field is $(2r+1)^2$ block tokens.
\end{itemize}

This parameter matters only for \texttt{conditioner="transformer"}. For
\texttt{conditioner="mlp"}, it is ignored.

\subsection{$Z_2$ Symmetry}

The code supports
\[
\texttt{parity}\in\{\texttt{none},\texttt{sym}\}.
\]

For \texttt{parity="sym"}, the scale and shift channels are projected to be
even and odd, respectively:
\begin{align}
s(u) &= \frac12 \bigl(\hat s(u) + \hat s(-u)\bigr), \\
t(u) &= \frac12 \bigl(\hat t(u) - \hat t(-u)\bigr).
\end{align}
This enforces
\[
s(-u)=s(u),\qquad t(-u)=-t(u),
\]
and therefore the full map respects the old-model-style
$z\mapsto -z$, $\phi\mapsto -\phi$ symmetry.

\section{Terminal $2\times 2$ Flow}

At the final scale, the last $2\times 2$ block is packed into a $4$-vector and
mapped by an unconditional RealNVP:
\[
g_{\mathrm{term}}:\mathbb{R}^4\to \mathbb{R}^4.
\]
It uses the same parity option as the checkerboard conditioners.

\section{Exact Log Density}

The inverse map returns
\[
(z,\log |\det J^{-1}|).
\]
Hence
\begin{equation}
\log p_\theta(x) = \log p_Z(z) + \sum_k \log |\det J_k^{-1}| + \log |\det J_{\mathrm{term}}^{-1}|.
\end{equation}

Each red or black subupdate contributes a diagonal affine Jacobian on the
updated color only. The level log-Jacobian is the sum over all red/black
subupdates in all cycles at that level.

\section{Training Objective}

Training uses the standard flow objective
\begin{equation}
\mathcal L(\theta) =
\mathbb E_{z\sim p_Z}
\bigl[
\log p_\theta(g_\theta(z)) + S(g_\theta(z))
\bigr].
\end{equation}

For reweighting quality, the more important validation statistic is not the raw
loss but the centered width of
\[
\Delta S = \log p_\theta(\phi) + S(\phi).
\]
In practice the useful proxy is
\[
\sigma\bigl(\Delta S - \langle \Delta S\rangle\bigr),
\]
because reweighting factors are $w = e^{-\Delta S}$.

\section{CLI Reference: Training Script}

This section documents
\texttt{scripts/phi4/train\_rg\_checkerboard\_flow.py}.

\begin{longtable}{p{0.24\linewidth}p{0.70\linewidth}}
\toprule
Argument & Meaning \\
\midrule
\endhead
\texttt{--resume} & Resume training from an existing checkpoint. \\
\texttt{--save} & Output checkpoint path. \\
\texttt{--save-every} & Save an intermediate checkpoint every $N$ epochs. Zero means only save at the end. \\
\texttt{--validate} & Run the validation routine at the end of training. In the current checkerboard trainer this is a final-only validation, not every 100 epochs. \\
\texttt{--validate-batch} & Batch size used inside one validation draw. \\
\texttt{--validate-super} & Number of validation draws to concatenate before computing $\Delta S$ and ESS summaries. \\
\texttt{--L} & Lattice linear size; the lattice is $L\times L$. Must be a power of two. \\
\texttt{--lam} & Quartic coupling $\lambda$. \\
\texttt{--mass} & Quadratic mass parameter $m^2$. \\
\texttt{--width} & Transformer model width at each checkerboard conditioner. \\
\texttt{--n-cycles} & Number of red/black plus shifted-red/black cycles per non-terminal RG level. \\
\texttt{--conditioner} & Conditioner family: \texttt{transformer} or \texttt{mlp}. \\
\texttt{--transformer-layers} & Number of transformer blocks in each conditioner network. Used only for \texttt{conditioner="transformer"}. \\
\texttt{--n-heads} & Number of attention heads. Used only for \texttt{conditioner="transformer"}. The width must be divisible by this number. \\
\texttt{--attn-radius} & Local attention radius on the block lattice. Used only for \texttt{conditioner="transformer"}. Radius zero means on-site conditioning only. \\
\texttt{--mlp-dim} & Hidden width of the transformer MLP blocks. Used only for \texttt{conditioner="transformer"}. Zero means use $2\times$\texttt{width}. \\
\texttt{--terminal-n-layers} & Number of RealNVP layers in the terminal $2\times 2$ flow. \\
\texttt{--terminal-width} & Hidden width of the terminal RealNVP MLPs. Zero means reuse \texttt{width}. \\
\texttt{--output-init-scale} & Small nonzero scale used to initialize the output projections near the identity while still allowing gradients to flow. \\
\texttt{--parity} & Symmetry mode. \texttt{sym} enforces the $\phi\to-\phi$, $z\to-z$ symmetry; \texttt{none} does not. \\
\texttt{--rg-type} & RG blocking rule: \texttt{average} or \texttt{select}. \\
\texttt{--log-scale-clip} & Tanh clip for affine log-scales, used for numerical stability. \\
\texttt{--epochs} & Total target epoch count. On resume this is the final epoch number, not the number of extra epochs. \\
\texttt{--batch} & Training batch size. \\
\texttt{--lr} & Adam learning rate. \\
\texttt{--seed} & Random seed. \\
\bottomrule
\end{longtable}

\section{CLI Reference: Checker Script}

This section documents
\texttt{scripts/phi4/check\_rg\_checkerboard\_flow.py}.

\begin{longtable}{p{0.24\linewidth}p{0.70\linewidth}}
\toprule
Argument & Meaning \\
\midrule
\endhead
\texttt{--shape} & Lattice shape used for invertibility and timing tests. \\
\texttt{--jac-shape} & Small lattice shape used for the dense autodiff Jacobian check. \\
\texttt{--batch-size} & Batch size used in invertibility and timing tests. \\
\texttt{--width}, \texttt{--n-cycles}, \texttt{--conditioner}, \texttt{--transformer-layers}, \texttt{--n-heads}, \texttt{--attn-radius}, \texttt{--mlp-dim}, \texttt{--terminal-n-layers}, \texttt{--terminal-width}, \texttt{--output-init-scale}, \texttt{--parity}, \texttt{--rg-type}, \texttt{--log-scale-clip} & Same architectural meanings as in the training script. \\
\texttt{--lam}, \texttt{--mass} & Theory parameters used in the timing loss-gradient benchmark. \\
\texttt{--tests} & Comma-separated subset of \texttt{invertibility}, \texttt{jacobian}, \texttt{timing}. \\
\texttt{--invert-tol} & Absolute tolerance for the forward-then-inverse reconstruction test. \\
\texttt{--jac-tol} & Tolerance for the autodiff-vs-analytic logdet comparison. \\
\texttt{--jac-perturb-scale} & Small perturbation applied to the initialized weights before the Jacobian test so the logdet is not trivially zero. \\
\texttt{--timing-warmup} & Warmup iterations for each timed kernel. \\
\texttt{--timing-iters} & Number of timed iterations per kernel. \\
\texttt{--json-out} & Optional JSON output path for the checker summary. \\
\texttt{--selfcheck-fail} & Exit with nonzero status if any selected test fails. \\
\bottomrule
\end{longtable}

\section{Current Practical Interpretation}

Empirically, the checkerboard update pattern seems to do the important
neighbor-coupling work. Because of that, a cheap conditioner can already work
well:
\begin{itemize}
\item small width,
\item one transformer layer,
\item attention radius zero,
\item or even a pure pointwise MLP conditioner,
\item symmetry on.
\end{itemize}

This suggests that the structural coupling induced by red/black updates plus
the shifted reblocking is more important here than using a large attention
window.

\end{document}
